\noindent\textbf{BEMERKUNGEN. ZUR NETZAUSGLEICHUNG.}\hfill327

\begin{center}\textbf{BEMERKUNGEN.}\end{center}

\par\noindent\hspace*{\parindent}Von der Netzausgleichung ist nur eine Anzahl loser, nicht numerirter Blätter vorhanden.
Eine vollständige in alle Einzelheiten gehende Wiederherstellung an der Hand derselben
--- wenn sie überhaupt möglich ist --- würde mit umständlichen Rechnungen verknüpft sein.

\par\noindent\hspace*{\parindent}Direct entnommen sind diesen Blättern Art.~[2], die Normalgleichungen des Art.~[4],
die Seitengleichungen im Art.~[5] und die Tabelle unter [9], der nur die Bezeichnungen
der Columnen zugefügt sind; nach ihnen zusammengestellt sind Art.~[3], die Tabellen
der Art.~[7] und [8], sowie die Normalgleichungen des Art.~[6]. Die Notiz [1] ist
einem Beobachtungs- und Rechnungsheft zur hannoverschen Gradmessung aus dem Jahre 1825 entnommen.

\par\noindent\hspace*{\parindent}Der Gang des \textsc{Gauss}schen Ausgleichungsverfahrens ist folgender.
Nachdem aus den Winkelgleichungen, der Methode der kleinsten Quadrate entsprechend,
die Normalgleichungen, S.~303 und 304, hergestellt waren, hat \textsc{Gauss} diese
(auf indirectem Wege, vergl.~S.~325) aufgelöst und daraus die Verbesserungen berechnet,
die die Dreieckswinkelgleichungen \textls[120]{allein} erfüllen.
Mit den erstmalig verbesserten Winkelwerthen wurden die 12 Seitengleichungen
in der \textsc{Gauss} eigenthümlichen Weise aufgestellt, S.~304---311.
Durch die modificirte Seitengleichung erreicht \textsc{Gauss}, dass, wenn sie nicht
mit einer andern zusammen hängt, immer nur die Winkelsummen der angrenzenden,
nicht aber die der innern Dreiecke der Figur, auf die sich die Seitengleichung bezieht,
durch die Ausgleichung beeinflusst werden; vergl.~den Brief an \textsc{Gerling}
vom 19.~Januar 1840, S.~253. Aus den umgeformten Seitengleichungen wurden die
12 Normalgleichungen, S.~312, hergeleitet, durch deren Auflösung neue Werthe
der Verbesserungen der Richtungen erhalten werden, welche die Seitengleichungen
\textls[120]{allein} erfüllen. Nach dieser ersten Ausgleichung werden wieder
die neu entstandenen Widersprüche der Dreiecke ausgeglichen, mit den gefundenen
Verbesserungen von neuem die Seitengleichungen und die zugehörigen Normalgleichungen
aufgestellt, die sich jetzt von den vorigen nur in den constanten Gliedern unterscheiden.
Durch ihre Auflösung ergeben sich wieder neue Werthe der Richtungsverbesserungen.
In dieser Weise hat \textsc{Gauss} die Ausgleichung noch zweimal wiederholt.
Dies Verfahren führt, wie im Supplementum theor.~comb.~observ.~Art.~19 gezeigt ist,
zu denselben Werthen, wie die Ausgleichung sämmtlicher Winkel- und Seitengleichungen
in einem Gusse. Dass die Rechnung so bald stehende Resultate ergeben hat,
liegt nach \textsc{Gauss} an der von ihm gewählten Form der Seitengleichungen, vergl.~S.~251.

\par\noindent\hspace*{\parindent}Substituirte man die Verbesserungen des Art.~[7] in allen 51 Winkelgleichungen,
so würde, wie man aus der Zusammenstellung des Art.~[8] ersieht, bei 11 Gleichungen
der Fehler $0$, bei 18 Gleichungen der Fehler $0^{\prime\prime}{,}001$,
bei 18 Gleichungen der Fehler $0^{\prime\prime}{,}002$ und bei 4 Gleichungen
der Fehler $0^{\prime\prime}{,}003$ sein.

\par\noindent\hspace*{\parindent}Die von den Punkten 28 bis 33 ausgehenden Richtungen haben von \textsc{Gauss}
nochmals Correctionen erhalten, deren Herkunft nicht ersichtlich ist.
Wahrscheinlich rühren sie von einer nachträglichen Änderung der beiden beobachteten
Winkel in Langwarden um $0^{\prime\prime}{,}128$ und $0^{\prime\prime}{,}015$ her.
Jedenfalls sind die unter [7] mitgetheilten Verbesserungen, die zur Ableitung
der folgenden Tabellen benutzt wurden, diejenigen, welche sich aus seiner Ausgleichung ergeben haben.

\par\noindent\hspace*{\parindent}Aus den 51 Dreieckswidersprüchen $w_i$ erhält man für den mittlern
Richtungsfehler $m$ nach der Näherungsformel
$\displaystyle m^2=\frac{1}{6\mathbin{.}51}\sum_{1}^{51}w_i^2$:
\[
m=\pm\,0^{\prime\prime}{,}643.
\]

\clearpage
\noindent328\hfill\textbf{BEMERKUNGEN. ZUR NETZAUSGLEICHUNG.}
\par\medskip

\par\noindent\hspace*{\parindent}Bei dem Anschluss des Punktes Hohenhorn an das Dreiecksnetz, S.~316,
vermittelst der Figur Hohenhorn.~21.~15.~17.~18 sind die Seitenverhältnisse
und die Winkel der Gradmessung festgehalten worden. Die Dreiecksseite Hamburg--Hohenhorn
diente der hannoverschen Grad- und Landesvermessung als Basis; vergl.~Bestimmung
des Breitenunterschiedes etc., S.~48, sowie den Brief an \textsc{Schumacher}
vom 22.~December 1827, S.~281.

\par\noindent\hspace*{\parindent}Zu der Tabelle, Art.~[9], ist folgendes zu bemerken.
Die Darstellung einer Dreiecksseite $i.k$ in der Ebene bilde mit der Geraden
durch ihren Anfangs- und Endpunkt die Winkel $\psi_{i.k}$ und $\psi_{k.i}$.
Bezeichnet dann $T$ das astronomische Azimuth, $t$ das Azimuth auf dem Sphäroid,
$\Theta$ das Azimuth in plano und $\mathrm{c}$ die Meridianconvergenz,
so ist (vergl.~S.~201):
\[
\begin{array}{r@{\qquad}c}
1)&t_{i.k}=T_{i.k}+c_i\\[4pt]
2)&\Theta_{i.k}=t_{i.k}-\psi_{i.k},\qquad
\Theta_{k.i}=\Theta_{i.k}\pm180^\circ,\qquad
t_{k.i}=\Theta_{k.i}+\psi_{k.i}.
\end{array}
\]

\par\noindent\hspace*{\parindent}Die Berechnung von $\psi$ erfolgt nach der Formel des Art.~15, S.~159;
vergl.~auch S.~216/217. Diese Formel, wie die zur Reduction der sphäroidischen
Dreiecksseite auf die ebene, verlangt die Kenntniss angenäherter Coordinaten der
Dreieckspunkte, die man sich durch eine vorläufige Berechnung von geringer Schärfe verschaffen muss.

\par\noindent\hspace*{\parindent}Ist also in dem sphäroidischen Dreiecksnetz das astronomische Azimuth
\textls[120]{einer} Dreiecksseite bekannt, so kann man durch successive Anwendung
der Gleichungen (2) und mit Hülfe der Winkel des Netzes $t$ und $\Theta$
für alle Seiten berechnen. Nun ist aber für alle Punkte des Hauptmeridians
$\mathrm{c}=0$. Geht daher die Dreiecksseite $1\,.\,2$, deren Azimuth beobachtet ist,
von einem Punkte desselben aus, so ist
\[
\begin{array}{r@{\qquad}c}
3)&t_{1.2}=T_{1.2}.
\end{array}
\]

\par\noindent\hspace*{\parindent}Zur Orientirung des hannoverschen Dreieckssystems dient die Seite Göttingen,
Theodolithplatz 1823---Nördliches Meridianzeichen; ihr Azimuth ist nach Art.~[9]:
$T_{1.2}=t_{1.2}=5^{\prime\prime}{,}473$.
In einem Briefe an \textsc{Gerling} vom 26.~December 1823
(vergl.~auch \textsc{Gerling}, Beiträge zur Geographie Kurhessens etc., S.~69)
gibt \textsc{Gauss} dafür $5^{\prime\prime}{,}471$ an.
Aus den Coordinaten des Theodolithplatzes von 1823 auf der Göttinger Sternwarte:
$x=-5{,}507\,\mathrm{m}$, $y=0$, und des nördlichen Meridianzeichens:
$x=-5019{,}756\,\mathrm{m}$, $y=-0{,}133\,\mathrm{m}$, Band~IV,
S.~416/417, folgt ebenfalls das Azimuth $=5^{\prime\prime}{,}471$.

\par\noindent\hspace*{\parindent}Sind nun die ebenen Azimuthe berechnet, wie es in der Tabelle unter [9]
geschehen ist, so kann man aus ihnen die Winkel der ebenen Dreiecke zusammenstellen
und alsdann, wenn die lineare Länge einer Dreiecksseite bekannt ist,
die Längen aller übrigen Dreiecksseiten berechnen. Man hat demnach nur
\textls[120]{einmal} nöthig, von einer sphäroidischen Dreiecksseite auf eine ebene
zu reduciren. Diese Übertragung geschieht nach der auf S.~215 gegebenen Formel;
wendet man sie auf alle sphäroidischen Dreiecksseiten an, so gelangt man
gleichfalls zu den ebenen Dreiecksseiten.

\par\noindent\hspace*{\parindent}Wenn aber die Seiten der ebenen Dreiecke und ihre Azimuthe bekannt sind,
so lassen sich aus ihnen leicht successive die ebenen rechtwinkligen Coordinaten
der Dreieckspunkte berechnen.

\par\noindent\hspace*{\parindent}Die auf S.~299 gegebene Karte des der Ausgleichung unterworfenen Dreiecksnetzes
ist die Copie einer von \textsc{Gauss} dem hannoverschen Cabinetsministerium
eingereichten Zeichnung. Aus Versehen sind im Original die Richtungen
Wilsede-Brüttendorf und Timpenberg-Hamburg fortgelassen; dagegen enthält dasselbe
noch den Anschluss der Punkte Hohenhorn und Lauenburg, sowie den der Punkte Wangeroog und Neuwerk.

\par\noindent\hspace*{\parindent}Über die Ausgleichung der in den ersten Jahren der Gradmessung beobachteten
Dreiecke sind die später folgenden Briefe an \textsc{Bessel} vom 5.~November 1823
und an \textsc{Bohnenberger} vom 16.~November 1823 zu vergleichen.
\hfill\textsc{Krüger.}

\begin{center}\rule{.24\textwidth}{.4pt}\end{center}

